\begin{frame}[plain]\FrameAnchor{V26-title}
\centering
{\Large\bfseries\color{courseblue}第 26 卷\par}
\vspace{0.35cm}
{\LARGE\bfseries 非微扰场论：对称破缺、有效理论、拓扑与禁闭\par}
\vspace{0.35cm}
{\large 建立微扰展开之外的真空、低能自由度和规范拓扑物理图像。\par}
\vfill
\CourseID{Tbl}{26.00}\quad 5 个学习单元，固定编号不随合订方式改变
\end{frame}

\begin{frame}[t]{第 26 卷路线图}\FrameAnchor{V26-route}
\small
\begin{tabularx}{\linewidth}{p{0.14\linewidth}Y}
\toprule 编号 & 学习单元\\\midrule
26.01 & 自发对称破缺与 Goldstone 定理\\
26.02 & Higgs 机制与规范玻色子质量\\
26.03 & 有效场论、OPE、matching 与系统幂计数\\
26.04 & 拓扑桥梁：S\ensuremath{^{3}}\ensuremath{\infty}\ensuremath{\rightarrow}SU(N)、\ensuremath{\pi}3=Z 与 instanton\\
26.05 & Wilson 圈、Polyakov loop 中心对称与禁闭边界\\
\bottomrule\end{tabularx}
\vfill
\SourceLine{BOOK-QFT；BOOK-LQCD；P01；P05；P07}
\end{frame}

\begin{frame}[t]{先修诊断与本卷出口}\FrameAnchor{V26-diagnostic}
\begin{columns}[T,onlytextwidth]
\column{0.48\textwidth}\begin{block}{开始前应能回答}
\small\begin{itemize}
\item V24
\item V25
\end{itemize}\end{block}
\column{0.48\textwidth}\begin{block}{学完后应能独立完成}
\small\begin{itemize}
\item 能应用 Goldstone 与 EFT 幂计数
\item 能解释 instanton/\allowbreak{}\ensuremath{\theta} 真空
\item 能用 Wilson 圈和中心对称讨论禁闭
\end{itemize}\end{block}
\end{columns}
\vfill\scriptsize 若先修项答不出，回到相应前卷；不要以术语熟悉代替可计算能力。
\end{frame}

\begin{frame}[t]{定量图：Mexican-hat 势的径向截面}\FrameAnchor{V26-chart}
\centering\begin{tikzpicture}
\begin{axis}[width=0.88\linewidth,height=0.63\textheight,xmin=-1.8,xmax=1.8,ymin=-0.3,ymax=2.0,xlabel={$\phi/v$},ylabel={$V/(\lambda v^4)$},grid=major,grid style={courseline!55},legend style={font=\scriptsize,draw=none,fill=white},tick label style={font=\scriptsize},label style={font=\small}]
\addplot[very thick,courseblue,domain=-1.8:1.8,samples=160] {(x^2-1)^2/4};
\addlegendentry{V=(phi\ensuremath{^{2}}-1)\ensuremath{^{2}}/\allowbreak{}4}
\end{axis}\end{tikzpicture}
\par\scriptsize\FigureID{26.00}：一分量截面；完整连续简并真空位于场空间圆/\allowbreak{}球面。
\SourceLine{自发对称破缺}
\end{frame}

\begin{frame}[t]{自发对称破缺与 Goldstone 定理：问题与图像}\FrameAnchor{26.01-1}
\footnotesize
\begin{block}{本节问题}方程对称而真空不对称时会出现什么低能模式？\end{block}
\PhysicalFlow{对称势 V}{选定非不变真空}{沿真空流形的无质量方向}{26.01}
\begin{block}{物理原则}\KnowledgeID{26.01}\quad 有限体积严格系综平均不会任选连续破缺方向；需外源、体积极限或相关函数诊断。\end{block}
\SourceLine{BOOK-QFT}
\end{frame}

\begin{frame}[t]{自发对称破缺与 Goldstone 定理：定义与主公式}\FrameAnchor{26.01-2}
\footnotesize\begin{tabularx}{\linewidth}{p{0.30\linewidth}Y}
\toprule 术语 & 本课程中的精确定义\\\midrule
\DefinitionID{26.01-9deb039680} 自发对称破缺 & 作用量对称但所选真空不在该对称下不变\\
\DefinitionID{26.01-bc939e0c8b} 序参量 & 区分不同相或真空的期望值\\
\DefinitionID{26.01-ba3b6e6ff3} Goldstone 模 & 每个破缺连续全局生成元对应的无质量低能激发\\
\bottomrule\end{tabularx}\vspace{0.15cm}
\begin{block}{\EquationID{26.01} 主公式}\centering\CourseDisplayEquation{V(\bm\phi)=\frac{\lambda}{4}(\bm\phi^2-v^2)^2,\qquad m_{\rm radial}^2=2\lambda v^2,\quad m_{\rm Goldstone}^2=0}\end{block}
势在 |\ensuremath{\phi}|=v 上简并；径向曲率非零，沿真空流形切向曲率为零。
\end{frame}

\begin{frame}[t]{自发对称破缺与 Goldstone 定理：推导与证据边界}\FrameAnchor{26.01-3}
\footnotesize
\DerivationID{26.01}\quad 由本节定义与前置结论逐步推出 \CourseRef{Eq}{26.01}：
\begin{enumerate}
\item 令 \ensuremath{\phi}=(v+h,\ensuremath{\pi}) 并展开 \ensuremath{\phi}\ensuremath{^{2}}-v\ensuremath{^{2}}=2vh+h\ensuremath{^{2}}+\ensuremath{\pi}\ensuremath{^{2}}。
\item 二次势只有 \ensuremath{\lambda}v\ensuremath{^{2}}h\ensuremath{^{2}}，对应 V含 m\ensuremath{^{2}}h\ensuremath{^{2}}/\allowbreak{}2，故 mh\ensuremath{^{2}}=2\ensuremath{\lambda}v\ensuremath{^{2}}。
\item \ensuremath{\pi} 不出现二次质量项，是沿简并真空流形的 Goldstone 方向。
\end{enumerate}\begin{block}{可执行检查}
\SympyID{26.01}\quad 自发对称破缺与 Goldstone 定理；2/2 项通过；SymPy exact。
\par\scriptsize 假设：O(2) Mexican-hat 势；在真空 (v,0) 展开
\end{block}
\BoundaryLine{有限体积自发破缺与量子修正未包含。}
\end{frame}

\begin{frame}[t]{自发对称破缺与 Goldstone 定理：计算算法}\FrameAnchor{26.01-4}
\footnotesize
\AlgorithmID{26.01}\begin{enumerate}
\item 写出对称群 G、真空稳定子 H 和破缺生成元数 dimG-dimH。
\item 在有限盒加入小外源选方向。
\item 测量序参量、纵/\allowbreak{}横向关联长度。
\item 按先 V\ensuremath{\rightarrow}\ensuremath{\infty} 后外源\ensuremath{\rightarrow}0 的顺序检查 Goldstone 极点。
\end{enumerate}\begin{columns}[T,onlytextwidth]
\column{0.56\textwidth}\begin{block}{停止条件与交叉检查}\scriptsize\begin{itemize}
\item[\CheckMark] v=0 时不破缺且径向/\allowbreak{}横向分类消失。
\item[\CheckMark] 显式小破缺会给 pseudo-Goldstone 小质量。
\end{itemize}\end{block}
\column{0.40\textwidth}\begin{block}{代码映射}
\scriptsize \texttt{课程内纸笔/\allowbreak{}符号计算练习}
\end{block}\end{columns}\end{frame}

\begin{frame}[t]{自发对称破缺与 Goldstone 定理：练习与完整解答}\FrameAnchor{26.01-5}
\footnotesize
\begin{block}{\ExerciseID{26.01}}O(2)\ensuremath{\rightarrow}O(1) 有几个 Goldstone 模？\end{block}
\begin{exampleblock}{\SolutionID{26.01}}dim O(2)-dim O(1)=1-0=1 个。\end{exampleblock}
\vfill
\scriptsize 回看：\CourseRef{Def}{26.01-9deb039680}；\CourseRef{Eq}{26.01}；\CourseRef{SYM}{26.01}；\CourseRef{Alg}{26.01}。
\SourceLine{BOOK-QFT}
\end{frame}

\begin{frame}[t]{Higgs 机制与规范玻色子质量：问题与图像}\FrameAnchor{26.02-1}
\footnotesize
\begin{block}{本节问题}Goldstone 模在局域对称性中去了哪里？\end{block}
\PhysicalFlow{复标量取得 v}{规范变换吸收相位}{纵向偏振补成有质量矢量}{26.02}
\begin{block}{物理原则}\KnowledgeID{26.02}\quad 规范对称不是物理上真正被破坏的冗余；可观测叙述应以规范不变量或固定规范后一致结果表达。\end{block}
\SourceLine{BOOK-QFT}
\end{frame}

\begin{frame}[t]{Higgs 机制与规范玻色子质量：定义与主公式}\FrameAnchor{26.02-2}
\footnotesize\begin{tabularx}{\linewidth}{p{0.30\linewidth}Y}
\toprule 术语 & 本课程中的精确定义\\\midrule
\DefinitionID{26.02-f592589402} Higgs 机制 & 规范场吸收 would-be Goldstone 并获得质量\\
\DefinitionID{26.02-26b2784a2c} unitary gauge & 把标量相位变换掉的规范\\
\DefinitionID{26.02-db05de47fa} 径向 Higgs 模 & 真空模长方向的物理标量激发\\
\bottomrule\end{tabularx}\vspace{0.15cm}
\begin{block}{\EquationID{26.02} 主公式}\centering\CourseDisplayEquation{|D_\mu\phi|^2\supset\frac12g^2v^2A_\mu A^\mu,\qquad m_A=gv}\end{block}
协变导数作用在非零真空上产生规范场二次项，即质量项。
\end{frame}

\begin{frame}[t]{Higgs 机制与规范玻色子质量：推导与证据边界}\FrameAnchor{26.02-3}
\footnotesize
\DerivationID{26.02}\quad 由本节定义与前置结论逐步推出 \CourseRef{Eq}{26.02}：
\begin{enumerate}
\item 写 \ensuremath{\phi}=(v+h)e\ensuremath{^{i\theta/v}}/\allowbreak{}\ensuremath{\sqrt{2}}、D\ensuremath{\mu}=\ensuremath{\partial}\ensuremath{\mu}-igA\ensuremath{\mu}。
\item 作规范变换令 \ensuremath{\theta}=0，D\ensuremath{\mu}\ensuremath{\phi} 包含 -igA\ensuremath{\mu}(v+h)/\allowbreak{}\ensuremath{\sqrt{2}}。
\item 模平方的 v\ensuremath{^{2}} 项为 g\ensuremath{^{2}}v\ensuremath{^{2}}A\ensuremath{^{2}}/\allowbreak{}2，故 mA\ensuremath{^{2}}=g\ensuremath{^{2}}v\ensuremath{^{2}}。
\end{enumerate}\begin{block}{可执行检查}
\SympyID{26.02}\quad Higgs 机制与规范玻色子质量；1/1 项通过；SymPy exact。
\par\scriptsize 假设：Abelian Higgs unitary-gauge 二次项代理
\end{block}
\BoundaryLine{不证明规范不变可观测谱或非 Abelian 混合。}
\end{frame}

\begin{frame}[t]{Higgs 机制与规范玻色子质量：计算算法}\FrameAnchor{26.02-4}
\footnotesize
\AlgorithmID{26.02}\begin{enumerate}
\item 从格点 action 构造规范—Higgs 系统并保持链接协变。
\item 测量规范不变复合算符而非裸 \ensuremath{\langle}\ensuremath{\phi}\ensuremath{\rangle}。
\item 用谱学提取标量和矢量质量。
\item 扫描体积、规范参数和相图，检查 Fradkin--Shenker 连续性等边界。
\end{enumerate}\begin{columns}[T,onlytextwidth]
\column{0.56\textwidth}\begin{block}{停止条件与交叉检查}\scriptsize\begin{itemize}
\item[\CheckMark] g\ensuremath{\rightarrow}0 时规范场质量消失，恢复全局 Goldstone 图像。
\item[\CheckMark] 自由度计数：无质量矢量 2 + Goldstone 1 \ensuremath{\rightarrow} 有质量矢量 3。
\end{itemize}\end{block}
\column{0.40\textwidth}\begin{block}{代码映射}
\scriptsize \texttt{课程内纸笔/\allowbreak{}符号计算练习}
\end{block}\end{columns}\end{frame}

\begin{frame}[t]{Higgs 机制与规范玻色子质量：练习与完整解答}\FrameAnchor{26.02-5}
\footnotesize
\begin{block}{\ExerciseID{26.02}}v 加倍且 g 不变，mA 如何变？\end{block}
\begin{exampleblock}{\SolutionID{26.02}}mA=gv，故加倍。\end{exampleblock}
\vfill
\scriptsize 回看：\CourseRef{Def}{26.02-f592589402}；\CourseRef{Eq}{26.02}；\CourseRef{SYM}{26.02}；\CourseRef{Alg}{26.02}。
\SourceLine{BOOK-QFT}
\end{frame}

\begin{frame}[t]{有效场论、OPE、matching 与系统幂计数：问题与图像}\FrameAnchor{26.03-1}
\footnotesize
\begin{block}{本节问题}局域算符为何能代替看不见的短程传播，系数又怎样由完整理论固定？\end{block}
\PhysicalFlow{短距离乘积展开成局域算符}{在共同低能外态上 matching}{按 E/\allowbreak{}\ensuremath{\Lambda} 和算符维数截断}{26.03}
\begin{block}{物理原则}\KnowledgeID{26.03}\quad matching 必须固定算符基、外态、尺度和方案；Wilson 系数与重整化算符各自依 \ensuremath{\mu}，乘积之和才与 \ensuremath{\mu} 无关。大系数、额外红外尺度、对称性选择规则和 loop 因子会修正仅靠工程维数的朴素估计。\end{block}
\SourceLine{BOOK-QFT；DOC-OPE}
\end{frame}

\begin{frame}[t]{有效场论、OPE、matching 与系统幂计数：定义与主公式}\FrameAnchor{26.03-2}
\footnotesize\begin{tabularx}{\linewidth}{p{0.30\linewidth}Y}
\toprule 术语 & 本课程中的精确定义\\\midrule
\DefinitionID{26.03-d8b3da0b01} 有效场论 & 只显式保留目标尺度以下自由度，并写出全部对称允许局域算符的有序展开\\
\DefinitionID{26.03-a5debd7e9a} OPE & 当 x\ensuremath{\rightarrow}0 时，在低能关联函数内部把 A(x)B(0) 渐近展开为 \ensuremath{\Sigma}n Cn(x,\ensuremath{\mu})On(0,\ensuremath{\mu})\\
\DefinitionID{26.03-b98477aed8} Wilson 系数 & 只含被积分掉短程尺度的系数，由完整理论与低能理论对同一振幅或矩阵元 matching 固定\\
\DefinitionID{26.03-935da18e6d} 幂计数 & 按 E/\allowbreak{}\ensuremath{\Lambda}、质量比和 loop 因子系统排序保留项与首个遗漏误差\\
\bottomrule\end{tabularx}\vspace{0.15cm}
\begin{block}{\EquationID{26.03} 主公式}\centering\CourseDisplayEquation{A(x)B(0)\underset{x\to0}{\sim}\sum_n C_n(x,\mu)O_n(0,\mu),\qquad \phi(x)\phi(0)\sim G(x)\mathbf1+{:}\phi^2(0){:}+x^\mu{:}(\partial_\mu\phi)\phi(0){:}+\cdots,\quad G(x)\sim\frac1{4\pi^2x^2};\qquad \mathcal L_{\rm EFT}=\mathcal L_{d\le4}+\sum_{d>4}\frac{C_dO_d}{\Lambda^{d-4}}}\end{block}
OPE 的符号 \ensuremath{\sim} 不是任意 x 的算符恒等式，而是 x 比所有外部长尺度都小时，在任意指定低能关联函数中按幂逐阶逼近。自由 Euclidean 标量在 m|x|\ensuremath{\ll}1 时的 Wick 收缩给 G(x)\ensuremath{\approx}1/\allowbreak{}(4\ensuremath{\pi}\ensuremath{^{2}}x\ensuremath{^{2}})，normal-ordered 剩余项再作 Taylor 展开便得到显示的短距例子。
\end{frame}

\begin{frame}[t]{有效场论、OPE、matching 与系统幂计数：推导与证据边界}\FrameAnchor{26.03-3}
\footnotesize
\DerivationID{26.03}\quad 由本节定义与前置结论逐步推出 \CourseRef{Eq}{26.03}：
\begin{enumerate}
\item 对自由标量用 Wick 定理写 \ensuremath{\phi}(x)\ensuremath{\phi}(0)=G(x)1+:\ensuremath{\phi}(x)\ensuremath{\phi}(0):；再令 \ensuremath{\phi}(x)=\ensuremath{\phi}(0)+x\ensuremath{\mu}\ensuremath{\partial}\ensuremath{\mu}\ensuremath{\phi}(0)+O(x\ensuremath{^{2}})，得到上式 OPE。取真空期望值立即 match 出 C1=G(x)；取 normal-ordered 两粒子矩阵元则固定 C\ensuremath{_{:\phi^{2}:}}=1。
\item 作为 EFT matching，再设完整理论有质量 M=\ensuremath{\Lambda} 的重传播子。低动量 q\ensuremath{^{2}}\ensuremath{\ll}M\ensuremath{^{2}} 时 1/\allowbreak{}(M\ensuremath{^{2}}-q\ensuremath{^{2}})=M\ensuremath{^{-2}}+q\ensuremath{^{2}}M\ensuremath{^{-4}}+...；让完整理论的低能散射振幅与局域接触算符及导数算符振幅逐阶相等，就固定各 Wilson 系数。
\item 四维中维数 d 算符以 \ensuremath{\Lambda}\ensuremath{^{4-d}} 配平；若其矩阵元只含尺度 E，典型相对抑制为 (E/\allowbreak{}\ensuremath{\Lambda})\ensuremath{^{d-4}}。同一阶全部允许算符都要保留，不能只挑容易计算的一项。
\item 从 matching 尺度 \ensuremath{\Lambda} 到实验尺度 E 用 RG 演化 Cn；算符 mixing 使 C 成向量、演化成矩阵，而物理和 \ensuremath{\Sigma}Cn\ensuremath{\langle}On\ensuremath{\rangle} 的 \ensuremath{\mu} 依赖逐阶相消。
\end{enumerate}\begin{block}{可执行检查}
\SympyID{26.03}\quad 有效场论、OPE、matching 与系统幂计数；1/1 项通过；SymPy exact。
\par\scriptsize 假设：四维时空；O\_d 的工程量纲为 d
\end{block}
\BoundaryLine{只验证 \ensuremath{\Lambda}\ensuremath{^{4-d}} 配平；Wilson 系数大小与 EFT 幂计数需物理输入。}
\end{frame}

\begin{frame}[t]{有效场论、OPE、matching 与系统幂计数：计算算法}\FrameAnchor{26.03-4}
\footnotesize
\AlgorithmID{26.03}\begin{enumerate}
\item 明确低能自由度、对称性、算符基和小参数。
\item 选完整/\allowbreak{}EFT 都可计算且无额外红外歧义的外态或 Green 函数，在 \ensuremath{\mu}\ensuremath{\approx}\ensuremath{\Lambda} matching。
\item 按幂计数保留目标阶全部算符并 RG 到 E。
\item 以首个遗漏阶给截断先验，再用改变 matching 尺度、加入下一阶和不同过程验证覆盖。
\end{enumerate}\begin{columns}[T,onlytextwidth]
\column{0.56\textwidth}\begin{block}{停止条件与交叉检查}\scriptsize\begin{itemize}
\item[\CheckMark] x\ensuremath{\rightarrow}0 时单位算符系数 G(x) 最奇异，而更高维算符伴随更多 x 幂。
\item[\CheckMark] E/\allowbreak{}\ensuremath{\Lambda}\ensuremath{\rightarrow}0 时高维项消失；遗漏同阶算符会使误差估计失去系统性。
\end{itemize}\end{block}
\column{0.40\textwidth}\begin{block}{代码映射}
\scriptsize \texttt{课程内纸笔/\allowbreak{}符号计算练习}
\end{block}\end{columns}\end{frame}

\begin{frame}[t]{有效场论、OPE、matching 与系统幂计数：练习与完整解答}\FrameAnchor{26.03-5}
\footnotesize
\begin{block}{\ExerciseID{26.03}}把自由标量 OPE 写到 O(x)，并说明怎样分别固定 1 与 :\ensuremath{\phi}\ensuremath{^{2}}: 的系数。\end{block}
\begin{exampleblock}{\SolutionID{26.03}}\ensuremath{\phi}(x)\ensuremath{\phi}(0)=G(x)1+:\ensuremath{\phi}\ensuremath{^{2}}(0):+x\ensuremath{\mu}:(\ensuremath{\partial}\ensuremath{\mu}\ensuremath{\phi})\ensuremath{\phi}(0):+O(x\ensuremath{^{2}})。取真空期望值时 normal-ordered 项为零，故 C1=G(x)；再取去掉真空收缩的两粒子矩阵元，树级两边相等给 C\ensuremath{_{:\phi^{2}:}}=1。\end{exampleblock}
\vfill
\scriptsize 回看：\CourseRef{Def}{26.03-d8b3da0b01}；\CourseRef{Eq}{26.03}；\CourseRef{SYM}{26.03}；\CourseRef{Alg}{26.03}。
\SourceLine{BOOK-QFT；DOC-OPE}
\end{frame}

\begin{frame}[t]{拓扑桥梁：S\ensuremath{^{3}}\ensuremath{\infty}\ensuremath{\rightarrow}SU(N)、\ensuremath{\pi}3=Z 与 instanton：问题与图像}\FrameAnchor{26.04-1}
\footnotesize
\begin{block}{本节问题}为什么有限作用量把连续的规范场分成不能平滑互相变形的整数扇区？\end{block}
\PhysicalFlow{把无穷远方向组成边界三球面}{纯规范边界给映射 U:S\ensuremath{^{3}}\ensuremath{\rightarrow}SU(N)}{同伦类、Q 与作用量下界相连}{26.04}
\begin{block}{物理原则}\KnowledgeID{26.04}\quad 纯规范表达式和 \ensuremath{\nu} 的总符号随 D\ensuremath{\mu}、取向与迹约定可同时翻转，但 Q=\ensuremath{\nu} 的内部约定必须一致。有限格距裸 clover Q 不严格整数；flow、改进与连续极限是诊断，不能以四舍五入制造拓扑。\end{block}
\SourceLine{BOOK-QFT；BOOK-LQCD；P05；P07}
\end{frame}

\begin{frame}[t]{拓扑桥梁：S\ensuremath{^{3}}\ensuremath{\infty}\ensuremath{\rightarrow}SU(N)、\ensuremath{\pi}3=Z 与 instanton：定义与主公式}\FrameAnchor{26.04-2}
\footnotesize\begin{tabularx}{\linewidth}{p{0.30\linewidth}Y}
\toprule 术语 & 本课程中的精确定义\\\midrule
\DefinitionID{26.04-f661b692b4} 同伦 & 两张映射若能在始终连续且保持边界条件的过程中互相变形，就属于同一类\\
\DefinitionID{26.04-ec36e9559f} 三阶同伦群 & \ensuremath{\pi}3(G) 记录从 S\ensuremath{^{3}} 到 G 的映射同伦类；\ensuremath{\pi}3(SU(N))=Z 对所有 N\ensuremath{\ge}2 成立\\
\DefinitionID{26.04-dc7f615a51} 绕数 & 映射 U:S\ensuremath{^{3}}\ensuremath{\rightarrow}SU(N) 包裹群流形的带方向整数次数\\
\DefinitionID{26.04-ff01446dd0} instanton & Euclidean 自对偶、有限作用量且拓扑荷非零的经典规范场解\\
\bottomrule\end{tabularx}\vspace{0.15cm}
\begin{block}{\EquationID{26.04} 主公式}\centering\CourseDisplayEquation{\begin{gathered}F_{\mu\nu}\to0\ (r\to\infty)\Rightarrow A_\mu\to\frac{i}{g}U\partial_\mu U^{-1},\qquad U:S^3_\infty\to SU(N),\\ \nu[U]=\frac1{24\pi^2}\int_{S^3_\infty}\operatorname{tr}(U^{-1}dU)^3\in\pi_3(SU(N))=\mathbb Z,\\ Q=\frac{g^2}{32\pi^2}\int d^4x\,F^a_{\mu\nu}\widetilde F^a_{\mu\nu}=\nu[U],\qquad S\ge\frac{8\pi^2}{g^2}|Q|.\end{gathered}}\end{block}
Euclidean R\ensuremath{^{4}} 中很大半径的边界只记录方向，拓扑上是 S\ensuremath{^{3}}\ensuremath{\infty}。有限作用量要求 F 衰减；边界上 F=0 的联络局部是纯规范，所以每个构型给一张 U:S\ensuremath{^{3}}\ensuremath{\infty}\ensuremath{\rightarrow}SU(N)。\ensuremath{\pi}3=Z 意味着其连续变形类由整数绕数标记。
\end{frame}

\begin{frame}[t]{拓扑桥梁：S\ensuremath{^{3}}\ensuremath{\infty}\ensuremath{\rightarrow}SU(N)、\ensuremath{\pi}3=Z 与 instanton：推导与证据边界}\FrameAnchor{26.04-3}
\footnotesize
\DerivationID{26.04}\quad 由本节定义与前置结论逐步推出 \CourseRef{Eq}{26.04}：
\begin{enumerate}
\item 有限作用量 S=(1/\allowbreak{}4)\ensuremath{\int}F\ensuremath{^{2}} 要求体积因子 r\ensuremath{^{3}} 增长时 F 足够快趋零；在无穷远球面上 F=0，因而 A 可写成 U 对零联络的规范变换。
\item 连续改变有限作用量构型会连续改变 U。不同整数绕数的 U 不可能在 S\ensuremath{^{3}}\ensuremath{\rightarrow}SU(N) 内连续互变，除非中途破坏边界条件或出现奇点；这就是拓扑扇区。
\item 由 F Ftilde 是 Chern--Simons current 的散度，四维积分化为 S\ensuremath{^{3}}\ensuremath{\infty} 边界积分，恰得到 \ensuremath{\nu}[U]，所以光滑有限作用量场的 Q 为整数。
\item 用 0\ensuremath{\le}\ensuremath{\int}(F\ensuremath{\mp}Ftilde)\ensuremath{^{2}}、Euclidean 中 Ftilde\ensuremath{^{2}}=F\ensuremath{^{2}} 平方完成，得到 S\ensuremath{\ge}8\ensuremath{\pi}\ensuremath{^{2}}|Q|/\allowbreak{}g\ensuremath{^{2}}；F=\ensuremath{\pm}Ftilde 时饱和，Q=\ensuremath{\pm}1 的解是基本 instanton/\allowbreak{}anti-instanton。
\end{enumerate}\begin{block}{可执行检查}
\SympyID{26.04}\quad 拓扑桥梁：S\ensuremath{^{3}}\ensuremath{\infty}\ensuremath{\rightarrow}SU(N)、\ensuremath{\pi}3=Z 与 instanton；3/3 项通过；SymPy exact。
\par\scriptsize 假设：S3\_infinity 以 x0\ensuremath{^{2}}+x1\ensuremath{^{2}}+x2\ensuremath{^{2}}+x3\ensuremath{^{2}}=1 表示；SU(2) 基本表示 U=x0 I+i x·sigma
\end{block}
\BoundaryLine{精确验证紧化空间的显式映射落在 SU(2)；pi\_3(SU(N))=Z、绕数积分整数性与 instanton 作用量界仍由拓扑和场论证明。}
\end{frame}

\begin{frame}[t]{拓扑桥梁：S\ensuremath{^{3}}\ensuremath{\infty}\ensuremath{\rightarrow}SU(N)、\ensuremath{\pi}3=Z 与 instanton：计算算法}\FrameAnchor{26.04-4}
\footnotesize
\AlgorithmID{26.04}\begin{enumerate}
\item 先在离散三球参数化一个已知 SU(2) 映射，数值积分 \ensuremath{\nu} 并随网格加密趋近整数。
\item 在多个 flow time 测 clover/\allowbreak{}改进 Q，同时检查作用量下界与整数平台。
\item 测 Q 的 Monte Carlo 时间序列、自相关和扇区隧穿率，诊断拓扑冻结。
\item 做多格距连续外推和体积检查；\ensuremath{\theta} 项则按 exp(i\ensuremath{\theta}Q) 加权各整数扇区。
\end{enumerate}\begin{columns}[T,onlytextwidth]
\column{0.56\textwidth}\begin{block}{停止条件与交叉检查}\scriptsize\begin{itemize}
\item[\CheckMark] N=1 的 SU(1) 平凡，\ensuremath{\pi}3=Z 的结论从 N\ensuremath{\ge}2 开始。
\item[\CheckMark] Q=0 只给 S\ensuremath{\ge}0；非零 Q 的最小作用量不能连续降到真空。
\end{itemize}\end{block}
\column{0.40\textwidth}\begin{block}{代码映射}
\scriptsize \texttt{课程内纸笔/\allowbreak{}符号计算练习}
\end{block}\end{columns}\end{frame}

\begin{frame}[t]{拓扑桥梁：S\ensuremath{^{3}}\ensuremath{\infty}\ensuremath{\rightarrow}SU(N)、\ensuremath{\pi}3=Z 与 instanton：练习与完整解答}\FrameAnchor{26.04-5}
\footnotesize
\begin{block}{\ExerciseID{26.04}}有限作用量在无穷远定义什么映射？其同伦分类和 Q=1 最小作用量分别是什么？\end{block}
\begin{exampleblock}{\SolutionID{26.04}}定义 U:S\ensuremath{^{3}}\ensuremath{\infty}\ensuremath{\rightarrow}SU(N)（N\ensuremath{\ge}2）；\ensuremath{\pi}3(SU(N))=Z，绕数等于拓扑荷 Q。Q=1 的 Bogomolny 下界为 8\ensuremath{\pi}\ensuremath{^{2}}/\allowbreak{}g\ensuremath{^{2}}，由自对偶 instanton 饱和。\end{exampleblock}
\vfill
\scriptsize 回看：\CourseRef{Def}{26.04-f661b692b4}；\CourseRef{Eq}{26.04}；\CourseRef{SYM}{26.04}；\CourseRef{Alg}{26.04}。
\SourceLine{BOOK-QFT；BOOK-LQCD；P05；P07}
\end{frame}

\begin{frame}[t]{Wilson 圈、Polyakov loop 中心对称与禁闭边界：问题与图像}\FrameAnchor{26.05-1}
\footnotesize
\begin{block}{本节问题}零温面积律与有限温度 ZN 对称分别怎样工作，动力学夸克又改变哪一步？\end{block}
\PhysicalFlow{大矩形 Wilson 圈提取静态势}{单时间片中心变换检验纯规范作用量}{Polyakov loop 作为中心荷一的序参量}{26.05}
\begin{block}{物理原则}\KnowledgeID{26.05}\quad 无限体积纯规范理论中未破缺中心相通常有 \ensuremath{\langle}L\ensuremath{\rangle}=0，破缺相可有非零 sector 值，并与单个静态 fundamental 色荷自由能相连。fundamental 动力学夸克的闭合世界线同样绕时间并取得 z，故显式破坏中心；此时 \ensuremath{\langle}L\ensuremath{\rangle} 不再是严格序参量，零温弦还可因成对产生轻夸克而断裂。\end{block}
\SourceLine{P01；BOOK-LQCD；BOOK-QCD-LATTICE}
\end{frame}

\begin{frame}[t]{Wilson 圈、Polyakov loop 中心对称与禁闭边界：定义与主公式}\FrameAnchor{26.05-2}
\footnotesize\begin{tabularx}{\linewidth}{p{0.30\linewidth}Y}
\toprule 术语 & 本课程中的精确定义\\\midrule
\DefinitionID{26.05-1267b0046a} 面积律 & 大 Wilson 圈期望按包围面积指数衰减，并给长距离线性静态势\\
\DefinitionID{26.05-db68a94c15} 中心 ZN & SU(N) 中与所有群元对易的 z=e\ensuremath{^{2\pi{}ik/N}}I\\
\DefinitionID{26.05-00dd0416c5} Polyakov loop & 在固定空间点沿周期 Euclidean 时间绕一周的 temporal links 乘积之归一化迹\\
\DefinitionID{26.05-e252340aba} 显式破缺 & fundamental 动力学物质不在中心变换下不变，使原本的全局 ZN 不再是作用量精确对称\\
\bottomrule\end{tabularx}\vspace{0.15cm}
\begin{block}{\EquationID{26.05} 主公式}\centering\CourseDisplayEquation{\langle W(R,T)\rangle\sim e^{-\sigma RT-\mu(R+T)},\quad V(R)=-\lim_{T\to\infty}T^{-1}\ln W(R,T)\sim\sigma R;\qquad L(\bm x)=\frac1N\operatorname{Tr}\prod_{\tau=0}^{N_t-1}U_4(\bm x,\tau),\quad U_4(\bm x,\tau_0)\to zU_4(\bm x,\tau_0),\ z\in Z_N:\quad L\to zL}\end{block}
中心变换把同一时间片 \ensuremath{\tau}0 上所有 temporal links 乘 z。每个纯规范 plaquette 含该片一个 U4 与一个反向 U4dagger，z z* 相消，所以 Wilson gauge action 与测度不变；绕时间一次的 Polyakov loop 只取一个 z，因而带中心荷 1。
\end{frame}

\begin{frame}[t]{Wilson 圈、Polyakov loop 中心对称与禁闭边界：推导与证据边界}\FrameAnchor{26.05-3}
\footnotesize
\DerivationID{26.05}\quad 由本节定义与前置结论逐步推出 \CourseRef{Eq}{26.05}：
\begin{enumerate}
\item 转移矩阵给 W(R,T)=\ensuremath{\Sigma}n|Zn(R)|\ensuremath{^{2}}e\ensuremath{^{-En(R)T}}；T\ensuremath{\rightarrow}\ensuremath{\infty} 选出最低静态能。面积律 e\ensuremath{^{-\sigma{}RT}} 在 -lnW/\allowbreak{}T 中留下 \ensuremath{\sigma}R，周长项只给静态自能常数和 1/\allowbreak{}T 修正。
\item 对单时间片作 U4\ensuremath{\rightarrow}zU4。逐个检查 spatial plaquette 不含该链接、temporal plaquette 中 z 与 z* 成对，所以纯规范 action 不变。
\item Polyakov 路径恰穿过该片一次，乘积获得 z；其共轭获得 z*。在精确未破缺 ZN 对称的有限体积系综中 sector 平均为零，研究自发破缺需看分布、|L|/\allowbreak{}susceptibility 和体积极限。
\item 加入 fundamental quark 后 hopping/\allowbreak{}worldline 跨该片只取得单个 z 且无法在 matter 项内抵消；中心成为显式破缺。Wilson 面积律、有限温度中心对称和有夸克弦断裂必须作为适用域不同的陈述。
\end{enumerate}\begin{block}{可执行检查}
\SympyID{26.05}\quad Wilson 圈、Polyakov loop 中心对称与禁闭边界；4/4 项通过；SymPy exact。
\par\scriptsize 假设：面积律加周长律模型；R,T>0；纯 SU(N) 规范理论中某时间片链接乘 z in Z\_N
\end{block}
\BoundaryLine{验证 Wilson 面积律代理及 L\_P->zL\_P、局域 plaquette 不变；动力学基本表示夸克显式破坏中心对称，弦断裂和 crossover 需有限温度数据。}
\end{frame}

\begin{frame}[t]{Wilson 圈、Polyakov loop 中心对称与禁闭边界：计算算法}\FrameAnchor{26.05-4}
\footnotesize
\AlgorithmID{26.05}\begin{enumerate}
\item 构造 Wilson 圈矩阵，用多路径/\allowbreak{}smearing 和 GEVP 提取 V(R)，扫描 R、T 窗。
\item 逐配置计算 L、其复平面分布、susceptibility 和 Binder/\allowbreak{}体积行为，并按 SU(N) 中心扇区检查。
\item 在代码中实施单时间片 z 变换，要求所有 plaquette/\allowbreak{}action 在舍入误差内不变且每个 L 精确乘 z。
\item 含动力学夸克时加入静态—轻介子对算符验证弦断裂，并明确 Polyakov loop 只作 crossover 诊断而非精确序参量。
\end{enumerate}\begin{columns}[T,onlytextwidth]
\column{0.56\textwidth}\begin{block}{停止条件与交叉检查}\scriptsize\begin{itemize}
\item[\CheckMark] bT 等无关：这里的中心变换作用于整个时间片，不是普通周期规范变换。
\item[\CheckMark] 短距离势应转为 Coulomb/\allowbreak{}运行耦合行为，不能把纯线性式外推到 R=0。
\end{itemize}\end{block}
\column{0.40\textwidth}\begin{block}{代码映射}
\scriptsize \texttt{课程内纸笔/\allowbreak{}符号计算练习}
\end{block}\end{columns}\end{frame}

\begin{frame}[t]{Wilson 圈、Polyakov loop 中心对称与禁闭边界：练习与完整解答}\FrameAnchor{26.05-5}
\footnotesize
\begin{block}{\ExerciseID{26.05}}对纯 SU(N) 格点作用量作单时间片 U4\ensuremath{\rightarrow}zU4：plaquette 和 Polyakov loop 分别怎样变？加入 fundamental quark 后呢？\end{block}
\begin{exampleblock}{\SolutionID{26.05}}每个 plaquette 的 z 与 z* 相消，纯规范 action 不变；每个绕时一次的 L 获得 z，故 L\ensuremath{\rightarrow}zL。fundamental quark 项/\allowbreak{}世界线获得未抵消的 z，中心对称被显式破坏，L 不再是严格序参量。\end{exampleblock}
\vfill
\scriptsize 回看：\CourseRef{Def}{26.05-1267b0046a}；\CourseRef{Eq}{26.05}；\CourseRef{SYM}{26.05}；\CourseRef{Alg}{26.05}。
\SourceLine{P01；BOOK-LQCD；BOOK-QCD-LATTICE}
\end{frame}

\begin{frame}[t]{第 26 卷能力清单}\FrameAnchor{V26-outcomes}
\small\begin{itemize}
\item[\CheckMark] 能应用 Goldstone 与 EFT 幂计数
\item[\CheckMark] 能解释 instanton/\allowbreak{}\ensuremath{\theta} 真空
\item[\CheckMark] 能用 Wilson 圈和中心对称讨论禁闭
\end{itemize}\vfill\begin{block}{通过标准}
不以“看懂”为标准：应能从空白纸推导主公式、实现算法、解释检查失败意味着什么，并完成本卷练习。
\end{block}\end{frame}

\begin{frame}[t]{第 26 卷来源（1/2）}\FrameAnchor{V26-sources}
\scriptsize\begin{tabularx}{\linewidth}{p{0.17\linewidth}p{0.28\linewidth}Y}
\toprule ID & 来源 & 本卷用途\\\midrule
\SourceID{BOOK-QFT} & An Introduction to Quantum Field Theory（仓库转排本） & 量子场论、规范理论、QCD 和重整化的主教材交叉核验。\\
\SourceID{BOOK-LQCD} & Introduction to Lattice QCD /\allowbreak{} 格点 QCD 导论（仓库转排本） & 欧氏化、规范链接、费米子、连续极限和关联函数。\\
\SourceID{P01} & 夸克禁闭 & 论文图谱 P01 的完整原始 PDF；底稿 books/\allowbreak{}复用(Wilson74)；SHA-256 6eadfb8d3e3cd217…。\\
\SourceID{P05} & 平凡化映射威尔逊流与HMC & 论文图谱 P05 的完整原始 PDF；底稿 arXiv:0907.5491；SHA-256 3320b1d4caaa4e90…。\\
\SourceID{P07} & 威尔逊流的性质与应用 & 论文图谱 P07 的完整原始 PDF；底稿 arXiv:1006.4518；SHA-256 63426411712d8b00…。\\
\bottomrule\end{tabularx}\end{frame}

\begin{frame}[t]{第 26 卷来源（2/2）}\FrameAnchor{V26-sources-2}
\scriptsize\begin{tabularx}{\linewidth}{p{0.17\linewidth}p{0.28\linewidth}Y}
\toprule ID & 来源 & 本卷用途\\\midrule
\SourceID{DOC-OPE} & 格点 QCD 中的 OPE 算符 & 局域矩、twist 与算符基底。\\
\SourceID{BOOK-QCD-LATTICE} & Quantum Chromodynamics on the Lattice（仓库转排本） & 格点 QCD 形式体系、算法和系统误差的深入交叉核验。\\
\bottomrule\end{tabularx}\end{frame}

